\subsection{Fundamento y límites de los escenarios y reservas}
La estimación conserva tres duraciones por actividad y contrasta sus efectos con red y recursos, siguiendo la separación entre duración, riesgo y calendario del manual de programación de NASA.~\parencite[secc.~6.3.1 y 7.4, pp.~38--40 y 59--60]{nasa-cronograma} Para la ponderación PERT se aplica $t_E=(t_O+4t_M+t_P)/6$ y $\sigma=(t_P-t_O)/6$ como sensibilidad de diseño. El escenario de tres valores conserva $t_O=t_M=d$ y $t_P=1,6d$ para ingeniería. La igualdad del optimista y probable representa el mínimo técnico del diseño inicial; el pesimista ensaya una recuperación adicional de hasta $0,6d$. Con ello $t_E=1,1d$, $\sigma=0,1d$ y $\sigma^2=0,01d^2$. Los ensayos conservan $(10,10,16)$, $t_E=11$ y $\sigma=1$ día. Revisiones contractuales, turnos y períodos de observación naturales no se ajustan con el mismo multiplicador.

Estos extremos son hipótesis de sensibilidad, no mediciones, percentiles ni una distribución calibrada. El registro \path{componentes/riesgos_reservas_planificacion.csv} identifica evento, paquetes, responsable, disparador, respuesta y límite de cada reserva. La recuperación de ingeniería deberá respaldarse por corrección y comprobación identificadas; si demanda más HH, se actualizarán paquetes y capacidad. La nivelación previa de $1,1d$ con las mismas HH comprueba disponibilidad temporal a menor intensidad; no financia automáticamente un 10\,\% de retrabajo ni demuestra probabilidad de entrega.

Los diez días hábiles marinos son una reserva de contingencia propuesta para una campaña con dos ventanas alternativas y cancelaciones acumuladas de hasta diez días en el escenario. No se atribuye esa duración a un pronóstico acreditado. Su uso requiere aviso o bloqueo registrado, decisión de Operación y autorización de la ventana alternativa. H3 incorpora su reserva en la cadena; H5 y H10 la preservan antes de habilitar su etapa. No se vuelve a contar como revisión, subsanación ni holgura libre de los cortes finales. Se registrarán duración disponible, consumida y restante; excederla exige escalar y recalcular dentro de los hitos contractuales.

El registro PL-R01--PL-R05 conserva evento, paquete, responsable, disparador y límite; T-16 enlaza sus tratamientos con los riesgos de marejada, integridad, contraparte y capacidad. PL-R01/R-01 consume una sola reserva por campaña y PL-R04 conserva la puerta de recepción E2 sin holgura. La actualización de un tratamiento se propagará a red y asignaciones mediante el control de cambios de SD6. La reserva de gestión se mantiene fuera de la línea base y su utilización requiere decisión de dirección y control de cambios; ninguna reserva autoriza desplazar el plazo contractual. La estimación funcional se concilia mediante el esfuerzo seleccionado y los lotes adicionales. La sensibilidad del ambiente y la firma externa de recepción E2 requieren control propio; no se cubren mediante una reserva meteorológica. No se afirma una confianza normal de dos sigmas sin sustentar distribuciones, correlaciones y efecto de las rutas que cambian bajo recursos.
